\begin{table}[tp]
\centering
\caption{Third-derivative-smooth nonconvex optimization: \href{https://github.com/k3nfalt/longtails/blob/main/LongTails/Setups/ThirdDerivativeSmoothNonconvex.lean}{ThirdDerivativeSmoothNonconvex}.}\label{tab:ThirdDerivativeSmoothNonconvex}
\parbox{\linewidth}{\footnotesize On a finite-dimensional real inner-product space, $f\in C^3$ has $L_3$-Lipschitz third derivative, $f(x_0)-\inf f\le\Delta$, and $\|\nabla^2f(x_0)\|\le H_0$, where $L_3,\Delta,\varepsilon>0$ and $H_0\ge0$. Exact value/gradient and full-Hessian oracles are available. The target is $\|\nabla f(x_{\rm out})\|\le\varepsilon$ almost surely; \textbf{(1)} guarantees it for every seed. No global gradient- or Hessian-Lipschitz bound or algorithm restriction is imposed.}\par\smallskip
\begin{threeparttable}
\footnotesize
\setlength{\tabcolsep}{4pt}\renewcommand{\arraystretch}{1.25}
\begin{tabularx}{\linewidth}{@{}>{\raggedright\arraybackslash}p{0.20\linewidth}>{\raggedright\arraybackslash}p{0.32\linewidth}>{\raggedright\arraybackslash}p{0.25\linewidth}>{\raggedright\arraybackslash}X@{}}
\toprule
\textbf{Result} & \textbf{Oracle complexity} & \textbf{Novelty} & \textbf{Limitations} \\
\midrule
\textbf{(1)} \href{https://github.com/k3nfalt/longtails/blob/main/LongTails/Theorems/LazyQuarticHessianTradeoff.lean}{Lazy\allowbreak Quartic\allowbreak Hessian\allowbreak Tradeoff} \textbf{(Ours)} $(\star\star\star)$ \href{https://github.com/k3nfalt/longtails/blob/main/scripts/reviews/LazyQuarticHessianTradeoff.md}{(review)}
& $\mathcal{O}\!\left((1+S^{5/4})\ell+m(1+S^{3/4})\right)$ total; $\mathcal{O}(1+S^{3/4})$ Hessian calls.
& & \\
\midrule[1.2pt]
\textbf{(2)} \citep{tdsDoikovGrapiglia2025v1}, Thm.~3.5, Cor.~3.6 ($p=3$)
& $\mathcal{O}\!\left(m\left[d^{1/3}\Delta\bar L^{1/3}\varepsilon^{-4/3}\right.\right.$\newline
$\left.\left.{}+d\{1+\log_+(L_3/L_{\rm init})\}\right]\right)$ total, counting trials and finite differences; $L_{\rm init}>0$.
& \textbf{(1)}: Dimension-free; for $S\ge1$, $L_{\rm init}\asymp L_3$, Hessians $S^{3/4}$ instead of $d^{1/3}S+d$; total gain when $md^{1/3}\gg S^{1/4}\ell$.
\newline Example: $m=d$, $H_0\le\beta$, $L_{\rm init}\asymp L_3$: $\mathcal{O}(d^{4/3}S+d^2)$ becomes $\mathcal{O}(S^{5/4}[1+\log S]+dS^{3/4})$.
& \textbf{(1)} needs $L_3,\Delta,H_0$; \textbf{(2)} adapts to $L_3$. For $m=d=1$, $L_{\rm init}\asymp L_3$, $S\ge1$: $\mathcal{O}(S)$ becomes $\mathcal{O}(S^{5/4}\ell)$. \\
\midrule
\textbf{(3)} \citep{tdsHinder2018v2}, first-order variant after (8); globally $L_1$-Lipschitz gradient
& $\mathcal{O}\!\left(d(1+S)\left[1+\log_+\!\left(\frac{L_1^3}{L_3\varepsilon^2}\right)\right]\right)$ expected value/gradient calls; no Hessians.
& \textbf{(1)}: Scope extends beyond global $L_1$; dimension-free. On the overlap, for $S\ge1$, gains when $d\gg S^{1/4}$ and $m\ll dS^{1/4}$, up to logarithms.
& \textbf{(1)} needs Hessians; larger cost when $d\lesssim S^{1/4}$, up to logarithms. \\
\midrule
\textbf{(4)} \citep{tdsCarmonEtAl2017}, Thm.~2; globally $L_1$-Lipschitz gradient
& $\mathcal{O}\!\left(\frac{\Delta L_1^{1/2}L_3^{1/6}}{\varepsilon^{5/3}}\left[1+\log_+\!\left(\frac{L_1\Delta}{\varepsilon^2}\right)\right]\right)$ value/gradient calls; no Hessians.
\newline Valid for $\varepsilon^{2/3}\le\delta^{1/2}L_3^{1/6}$ and $8L_3^{1/3}\varepsilon^{2/3}\le L_1$.
& \textbf{(1)}: Scope extends beyond global $L_1$, including quartic least squares. On the overlap, replaces $\sqrt{L_1}$ by $(L_3\Delta)^{1/4}$; total gain for $L_1\gg\beta$ and $m\ll\Delta^{1/4}L_1^{1/2}L_3^{-1/12}\varepsilon^{-2/3}$, up to logarithms.
& \textbf{(1)} needs Hessians and $H_0$; larger gradient prefactor when $L_1<\beta$ within \textbf{(4)}'s admissible accuracy range. \\
\bottomrule
\end{tabularx}
\begin{tablenotes}[flushleft]\footnotesize
\item[] \textbf{Abbr:} $d=\dim E$; $m\ge1$ is the cost of one full Hessian in value/gradient-call units; total cost is $G+mH$, with $G,H$ the respective call counts. $S=\Delta L_3^{1/3}\varepsilon^{-4/3}$; $\beta=\sqrt{L_3\Delta}$; $\ell=1+\log(1+H_0\Delta/\varepsilon^2+L_3\Delta^3/\varepsilon^4)$. $L_1$ is a global gradient-Lipschitz constant; $L_{\rm init}$ is \textbf{(2)}'s initial smoothness estimate and $\bar L=\max\{L_{\rm init},L_3\}$. $\delta=f(x_0)-\inf f\le\Delta$; $\log_+t=\max\{0,\log t\}$; $\asymp$ denotes equality up to universal factors.
\end{tablenotes}
\end{threeparttable}
\end{table}
\clearpage
